% Chapter3/chapter3.tex -- System Architecture and Design
% -----------------------------------------------------------------------------
% This chapter describes WHAT you built and WHY you designed it that way.
% The Implementation chapter (Chapter 4) then explains HOW you built it in code.
%
% Good chapter structure for an architecture chapter:
%   3.1 -- Top-level system overview (the big picture)
%   3.2 -- Hardware/software partitioning decisions
%   3.3 -- Individual component designs
%   3.4 -- Data flow and memory architecture
%   3.5 -- Quantisation scheme

\chapter{System Architecture and Design}
\label{ch:architecture}

\section{Top-Level System Overview}
\label{sec:system_overview}

The complete inference system consists of two main subsystems that run
concurrently on the ZCU102's heterogeneous architecture: a software host
application running on the ARM Cortex-A53 processing system (PS), and a
set of HLS-generated accelerator kernels executing in the Zynq UltraScale+
programmable logic (PL). The PS handles tokenisation, orchestrates the
overall autoregressive generation loop, and manages data transfers between
DDR memory and the PL accelerators via the AXI interconnect. The PL
accelerators perform the computationally intensive operations -- matrix
multiplication, attention, and element-wise transformations -- in a pipelined
hardware datapath.

% -- Suggested figure: system block diagram ------------------------------------
% \begin{figure}[htbp]
%   \centering
%   \includegraphics[width=0.9\linewidth]{Chapter3/media/system_block_diagram.png}
%   \caption{Top-level system architecture showing the PS--PL partitioning
%            on the ZCU102. The ARM Cortex-A53 runs the host application
%            while the FPGA fabric hosts the HLS accelerator kernels.}
%   \label{fig:system_block_diagram}
% \end{figure}

\lipsum[1] % Replace with your actual system overview

\section{Hardware/Software Partitioning}
\label{sec:hw_sw_partition}

The decision of which operations to accelerate in hardware and which to
run in software on the ARM core is driven by two criteria: computational
intensity (operations per byte of memory accessed) and implementation
complexity. Operations with low arithmetic intensity -- such as tokenisation,
sampling (temperature scaling and argmax over the vocabulary), and residual
stream additions -- are left on the CPU, where their low latency is not a
bottleneck. The computationally dominant operations -- the matrix--vector
multiplications in the attention and MLP blocks, which together account for
over 95\% of all arithmetic -- are offloaded to the PL.

\lipsum[2] % Replace with your justification of partitioning decisions

\section{Transformer Layer Architecture}
\label{sec:layer_architecture}

Each of TinyLlama's 22 transformer layers follows an identical structure,
which simplifies the hardware design considerably: a single universal
kernel design can be time-multiplexed across all layers rather than requiring
22 separate accelerators. Each layer consists of two sub-blocks in sequence:
the multi-head attention (MHA) block and the feed-forward network (FFN) block,
both preceded by RMSNorm layer normalisation.

\subsection{RMSNorm}
\label{subsec:rmsnorm}

Root Mean Square Layer Normalisation (RMSNorm) normalises the activation
vector $\mathbf{x} \in \mathbb{R}^{d_{\text{model}}}$ by its RMS magnitude:

\begin{equation}
  \text{RMSNorm}(\mathbf{x}) = \frac{\mathbf{x}}{\sqrt{\frac{1}{d}\sum_{i=1}^{d} x_i^2 + \epsilon}} \odot \mathbf{g}
  \label{eq:rmsnorm}
\end{equation}

where $\mathbf{g} \in \mathbb{R}^d$ is a learned element-wise scale vector
and $\epsilon$ is a small constant for numerical stability. Unlike the more
common LayerNorm, RMSNorm omits the mean-centring step, reducing the number
of passes required over the input vector from two to one -- a meaningful
simplification for hardware implementation.

\lipsum[3] % Replace with your RMSNorm hardware design description

\subsection{Multi-Head Attention}
\label{subsec:attention}

\lipsum[4] % Replace with your MHA design description

\subsection{SwiGLU Feed-Forward Network}
\label{subsec:swiglu}

\lipsum[5] % Replace with your FFN design description

\section{Memory Architecture and Data Flow}
\label{sec:memory_arch}

\lipsum[6] % Replace with your memory tiling, ping-pong buffering, etc.

\section{Quantisation Scheme}
\label{sec:quant_scheme}

All weight tensors are quantised to INT8 using symmetric per-channel
post-training quantisation (PTQ). For a weight tensor $W$, the quantised
representation $\hat{W}$ is computed as:

\begin{equation}
  \hat{W} = \text{clip}\!\left(\text{round}\!\left(\frac{W}{s}\right), -128, 127\right)
  \label{eq:quantisation}
\end{equation}

where the per-channel scale factor $s = \max(|W|) / 127$ is chosen such
that the maximum absolute value in the channel maps to the INT8 range
$[-128, 127]$ without overflow. Activations are quantised dynamically at
runtime using a running estimate of the activation range collected during
a calibration pass over representative inputs.

\lipsum[7] % Replace with your full quantisation analysis and justification
